% Transcription — fonds Alexandre Grothendieck, shelfmark 143, batch 2 (pages 21-33).
% Established from the facsimile archives/batches/143/batch-02.pdf, cut from a
% copy of 143.pdf supplied by hand (the archive host could not be reached from
% the session); PDF page k is the archivists' page 20+k, no cover sheet. Per
% the skill transcribe-grothendieck. The folder is thirty-three pages; this is
% the second and last batch. Batch 1 (pages 1-20) was being transcribed in
% parallel by another agent and did not exist when this pass was made, so the
% notation below is fixed from these pages alone, checked against
% transcripts/140-1 and transcripts/145 for the usage of the period.
%
% Pass: Opus 5.5 (claude-opus-5-5), 2026-09-24 — first pass, unchecked against the pages by a human.
%
% Pages skipped: 33, a computer listing (columns of decimal numbers, upside
% down) carrying only two short numerical sums in another context (« ΔU = ...
% m/s », « U_m = ... m/s »), nothing of the mathematics in hand. No page is
% summarised: there is no typescript in this batch. Page 32 is written, in
% landscape, over a printed IBM listing; the listing is ignored and his ink is
% transcribed.
%
% His own pagination: 11 on page 21, 12 on page 23, 13 on page 25, 14 on page
% 27, 15 on page 28, 16 on page 30 — recorded as \note{p. N de l'auteur}.
% Pages 22, 24, 26 and 29 carry no number of his (they read as the versos of
% the numbered sheets, and the text runs on across them); pages 31 and 32 carry
% none either. The only titles of his are « (9) Les schémas et Topos
% fondamentaux ... » (page 26), its sub-heading « (A) Topos modulaires ... »,
% and « Une conjecture arithmético-géométrique » (page 28, under a circled
% numeral read 9); the first heading (page 21) is the edition's and says so.
%
% What the batch is about. Pages 21-26 (his pp. 11-13) continue a discussion
% begun in batch 1 of how to encode, by groupoids, the action of a group Sigma
% (an extension of G by pi) on a system (T, pi): the « T-groupoïde » <T, pi>
% obtained by identifying the fundamental groups of its components with a
% fixed T, the functor from <T, pi> to the groupoid of universal coverings of
% Pi, and conversely a T-groupoid Lambda with a faithful functor to <pi> whose
% Isom sets pi_{t,t'} form, for each t, a « pseudo-partition » of pi — whence
% Lambda ≃ <T, pi>. He finds this « pas bien éclairant » and turns to a
% functor F : Rev(Pi) -> (Ens), E |-> T° ∧_pi E, commuting with arbitrary
% inductive limits and compatible with the operations of Sigma, then breaks
% off (« Arrêtons ... »). Pages 26-27 (his § 9, « Les schémas et Topos
% fondamentaux », A « Topos modulaires ») treat the modular topoi T_{0,nu} of
% genus-0 curves with nu marked points: for nu = 0, 1, 2 they are the
% classifying topoi B_PGL(2), B_Aff(1), B_Gm, with trivial pi_1 for nu <= 3;
% for nu >= 4, T_{0,nu} is the scheme of nu-3 distinct points of
% P^1 - {0,1,∞}, whose fundamental group is (the profinite completion of) a
% braid group Pi_{0,nu}, with S_nu x Gamma acting. Pages 28-30 (his pp. 15-16)
% state « Une conjecture arithmético-géométrique »: for the modular topos
% M_{g,nu} (an algebraic étendue except for (g,nu) = (0,0),(0,1),(0,2),(1,0))
% the natural homomorphisms pi_1(M_{g,nu;ξ}) -> Autext_lac°(pi_1(X̄ -
% {x_1,...,x_nu})) (and variants °° for the geometric part, and the
% Teichmüller group T_{g,nu}) should be isomorphisms; diagram (4) of two short
% exact sequences, the hoped-for description of Gamma as Autext_lac° / T^_{g,nu}
% with the example g = 0, nu = 3 recovering Autext_lac°(pi_1(P^1 - {0,1,∞})),
% sequence (5) with the quotient N_{g,nu} = M_{g,nu}/S_nu, a later NB on the
% automorphism group G_{g,nu} of the geometric modular topos, and the question
% whether Gamma -> Autext(T^_{g,nu}) is injective / bijective. Page 31 has one
% line: the question Autext(T_{g,nu}) = 1 ?. Page 32 (landscape, on a
% listing) draws two staggered diagrams of short exact sequences: the fibration
% pi_1(X̄_ξ - S_nu) -> pi_1(M̄_{g,nu+1}) -> pi_1(M̄_{g,nu}) over Gamma, and its
% profinite, « Autlac° / Autextlac° » counterpart with T^_{g,nu+1}, T^_{g,nu},
% labelled Sigma_{nu+1}, Sigma_nu, with scattered notes (M'_{g,nu}(k) ∋ ξ,
% T'_{g,nu} ∋ G_ξ, Sigma_{g,nu}).
%
% Notation fixed once for the batch, from his hand: \mathcal{T} for the set of
% objects (pp. 21-25), for the modular topoi \mathcal{T}_{0,\nu} (pp. 26-27)
% and for the Teichmüller groups \mathcal{T}_{g,\nu}, \widehat{\mathcal{T}}_{g,\nu}
% (pp. 28-32) — it is the same script letter on the page; M_{g,\nu}, \bar
% M_{g,\nu} (he writes M over a first T on page 28) and N_{g,\nu}; \Pi for the
% curve's group, \pi for its fundamental group, \widetilde{\Pi} for the universal
% covering; \mathfrak{S}_\nu; \Gamma (the Galois group, never named on these
% pages); « Autextlac » set \mathrm{Autext}_{\mathrm{lac}}, with his exponents
% ° and °° and the ! of page 32; \mathrm{Rev}(\Pi), (\mathrm{Ens}),
% (\pi\text{-}\mathrm{Ens}); B_{\mathrm{PGL}(2)} etc. for the classifying topoi.
%
% Hardest pages: 22 and 24 (fast prose, many connectives lost), 28 (the
% exponents and interlinear additions around (*), (**), (***) and a diagonal
% marginal note), 29 (two diagonal marginal notes, one block heavily scribbled
% out), and page 32 (layout of the staggered diagrams). Readings to check
% first: « pseudo-partition » (p. 24); « modulaires » in the heading of (A)
% (p. 26); the list B_PGL(2), B_Aff(1), B_Gm (p. 26); « tresses » (p. 27);
% « associé » and « étendue algébrique » (p. 28); the exponents ° / °° and the
% bar on X in (*)-(***) (p. 28); the vertical arrow label of (4) (p. 29).

\documentclass[11pt,a4paper]{article}
% Le préambule commun à toutes les transcriptions du fonds Grothendieck.
%
% Il tient en une idée : **l'incertitude se compose, elle ne se résout pas.**
% Une transcription qui lisse un mot douteux a détruit la seule chose pour
% laquelle on la faisait — un lecteur doit pouvoir savoir, à chaque endroit,
% ce qui était sur le feuillet et ce qui a été ajouté. Les six macros qui
% suivent sont donc l'appareil critique, et non de la décoration.
%
% Les mêmes commandes sont reconnues par `scripts/render.mjs`, qui produit la
% vue de lecture du site. Une macro ajoutée ici doit l'être aussi là-bas,
% faute de quoi le rendu échouera bruyamment — ce qui est voulu : un rendu
% silencieusement faux est pire qu'un rendu absent.

\ProvidesPackage{grothendieck}[2026/08/08 Transcriptions du fonds Grothendieck]

\usepackage{amsmath,amssymb,amsthm}
\usepackage{mathrsfs}
\usepackage{stmaryrd}
% Les diagrammes commutatifs sont partout dans ce fonds : c'est la langue
% même de la géométrie algébrique de Grothendieck, et une transcription qui
% les rendrait en prose ne transcrirait rien.
\usepackage{tikz-cd}
% Français par défaut — c'est la langue des feuillets — mais l'anglais est
% chargé aussi : la traduction et le résumé sont des documents anglais, et la
% typographie française (espaces avant les deux-points) y serait une faute.
% Ces documents basculent par \selectlanguage{english} en tête de corps.
\usepackage[english,french]{babel}
\usepackage{fontspec}
\usepackage{geometry}
\geometry{a4paper,margin=25mm,marginparwidth=28mm}
\usepackage{marginnote}
\usepackage{xcolor}
% ulem, not soul. `soul`'s \st and \ul are built on a character-by-character
% scan that XeLaTeX's font machinery breaks: the document fails with an
% undefined control sequence rather than degrading. ulem does the same two jobs
% and is Unicode-engine safe. `normalem` stops it hijacking \emph, which the
% transcriptions use for Grothendieck's own emphasis.
\usepackage[normalem]{ulem}
\usepackage{hyperref}
\hypersetup{colorlinks=true,linkcolor=black,urlcolor=black,pdfborder={0 0 0}}

% --- Métadonnées du lot -----------------------------------------------------
% Reprises telles quelles par le script de rendu, qui les affiche en tête de
% la vue de lecture. Elles fixent ce que le fichier prétend couvrir : sans
% elles, une transcription flotte sans point d'attache dans un fonds de seize
% mille feuillets.

\newcommand{\folder}[1]{\gdef\grothFolder{#1}}
\newcommand{\batch}[1]{\gdef\grothBatch{#1}}
\newcommand{\leaves}[2]{\gdef\grothFirst{#1}\gdef\grothLast{#2}}
\newcommand{\foldertitle}[1]{\gdef\grothFoldertitle{#1}}
\gdef\grothFolder{?}\gdef\grothBatch{?}\gdef\grothFirst{?}\gdef\grothLast{?}\gdef\grothFoldertitle{}

% --- Le résumé --------------------------------------------------------------
%
% Il ouvre la lecture modernisée, sous le titre « Résumé », et il est composé
% en petit. Ce n'est pas un ornement : c'est le seul passage du document qui
% ne soit pas des mathématiques, et un lecteur doit pouvoir le reconnaître
% avant de l'avoir lu — puis le sauter s'il connaît déjà le sujet, ou s'y
% arrêter s'il ne le connaît pas. Le filet qui le suit marque la reprise du
% corps.

\newenvironment{resume}
  {\small\setlength{\parskip}{0.45em}}
  {\par\medskip\hrule\medskip}

% --- Mots-clés ---------------------------------------------------------------
%
% En anglais, en clôture du résumé de la lecture modernisée : le vocabulaire
% moderne sous lequel chercher ce que les feuillets construisent. Le manifeste
% du site (`npm run manifest`) les extrait pour étiqueter la cote entière —
% cette ligne est la source unique des tags, il n'y a pas de fichier à côté.
\newcommand{\keywords}[1]{\par\smallskip\noindent
  {\footnotesize\textsc{Keywords} --- \textit{#1}}\par}

% --- L'appareil critique ----------------------------------------------------

% Le numéro de feuillet, en marge. C'est l'ancre qui permet de régler un
% désaccord sur une formule en regardant *un* feuillet plutôt qu'une chemise
% de six cents. Le site s'en sert aussi pour tourner les pages du fac-similé
% quand on fait défiler la transcription.
\newcommand{\leaf}[1]{%
  \marginnote{\footnotesize\color{gray!70}\textsf{#1}}%
  \label{leaf:#1}\ignorespaces}

% Illisible : on ne devine pas. Un mot inventé qui se lit comme les autres est
% le pire résultat possible.
\newcommand{\ill}{\textcolor{red!60!black}{[\,\ldots\,]}}

% Lecture douteuse : le mot est donné, et signalé comme incertain.
\newcommand{\uncertain}[1]{\uline{#1}}

% Ajout de l'éditeur — une lettre manquante, un mot sous-entendu.
\newcommand{\add}[1]{\textcolor{blue!50!black}{[#1]}}

% Biffé par Grothendieck lui-même. Ce qu'il a rayé fait partie du document :
% une rature dit souvent où la pensée a changé de direction.
\newcommand{\struck}[1]{\sout{#1}}

% Note du transcripteur, sur l'état matériel du feuillet.
\newcommand{\note}[1]{\par\smallskip{\small\color{gray!60!black}\textit{[#1]}}\par\smallskip}

% Note marginale de Grothendieck, distincte du corps du texte.
\newcommand{\marginal}[1]{\par\smallskip\noindent\hspace{1em}%
  {\small\color{gray!60!black}#1}\par\smallskip}

% --- Titre ------------------------------------------------------------------

\AtBeginDocument{%
  \begin{center}
    {\large\bfseries Fonds Alexandre Grothendieck --- cote n\textsuperscript{o}~\grothFolder}\\[2pt]
    {\itshape\grothFoldertitle}\\[4pt]
    {\small feuillets \grothFirst--\grothLast\ (lot \grothBatch)}\\[2pt]
    {\footnotesize\color{gray!60!black}Transcription établie d'après le fac-similé numérisé
     par l'Université de Montpellier.\\
     Les lectures incertaines sont soulignées, les illisibles notées [\,\ldots\,],
     les ajouts entre crochets.}
  \end{center}
  \vspace{1em}}

\endinput


\folder{143}
\batch{2}
\pages{21}{33}
\dating{[vers 1980-1981]}
\watermark{Édition de démonstration}
\foldertitle{Paradigmes des courbes algébriques (version provisoire) : notes manuscrites (s.d.)}

\begin{document}

\section*{Groupoïdes et opérations de $\Sigma$ (suite)}

\note{titre de l'éditeur : la page 21 continue sans titre un développement
commencé dans le premier lot}

\page{21}
\note{p. 11 de l'auteur}

\ldots\ un syst.\ transitif d'isom.\ entre les $\pi_t$ ($t \in \mathcal{T}$),
permettant de les identifier tous : \ill{} $T$. Ainsi $\langle \mathcal{T}, \pi
\rangle$ devient un $T$-groupoïde, i.e.\ un groupoïde (pas néc.\ connexe) où
les groupes fondamentaux des comp.\ connexes sont tous identifiés à $T$. De
plus, $\Sigma$ opère sur ce $T$-groupoïde (en opérant trivialement sur $T$) via
ses opérations sur le système $(\mathcal{T}, \pi)$ ($\pi$-\ill{}
\uncertain{prolongeant}), en opérant sur $\mathcal{T}$ par l'op.\ donnée et
sur $\pi$ via \ill{} \ill{} ext.\ de $\Sigma$, enfin sur $T$ via $G$
\uncertain{via} \ill{}\ldots

Si $\Sigma$ opère sur $\mathcal{T}$, \ill{} \ill{}, et \ill{} : $\pi$ \ill{},
on trouve un groupoïde $\langle \mathcal{T}, \pi \rangle$ sur lequel $\Sigma$
opère. Supposons \ill{} un groupoïde \ill{}, \ill{} donné \uncertain{par une}
famille de groupes \ill{} $(T_j)_{j \in J}$ indexée par $J = \mathcal{T}/\pi$,
sur laquelle opère $\Sigma$ via $\Sigma/\pi = G$. Ceci dit, on suppose donnée
une trivialisation de cette famille $(T_j)$ par $T \simeq T_j$, compatible
avec \ill{} \ill{}

\page{22}
$G$ --- lequel opère donc sur $T$. Ceci donne naissance automatiquement
\ill{} : \ill{} $\mathcal{T} \to$ \ill{}\ldots

Ceci dit, on a un hom.\ du groupoïde $\langle \mathcal{T}, \pi \rangle$ dans
le groupoïde des revêt.\ universels de $\Pi$, construits sur les \ill{} ;
explic.\ : \ill{} à valeurs $\widetilde{\Pi}$ sur ceux-ci, \ill{} dans l'ens.\
des \ill{} dans $\pi$ --- et pour $t, t' \in \mathcal{T}$, pour
\[
  \pi_{t,t'} \;\simeq\; \mathrm{Isom}_{\langle \mathcal{T}, \pi \rangle}(t, t')
\]
\note{l'isomorphisme avec $\mathrm{Isom}$ est ajouté au-dessus de la ligne et
rattaché par un trait}
l'ens.\ des $\lambda \in \pi$ tels que $\lambda t = t'$, \ill{} l'inclusion de
$\pi_{t t'}$ dans $\pi$. Le foncteur
\[
  \langle \mathcal{T}, \pi \rangle \longrightarrow
  \text{Cat des revêt.\ universels de } \Pi
\]
\ill{} ($\widetilde{\Pi}$ ; \uncertain{un seul objet} et groupe
d'\uncertain{isom.} $\pi$) \ldots

\ill{} $\Sigma$ sur $\Pi$ \ill{} donné sur $\pi$ \ldots\ compatibles \ill{}
\ill{} \ldots

\page{23}
\note{p. 12 de l'auteur}

Inversement, donnons-nous un $T$-groupoïde \struck{\ill{}} $\Lambda$, (un
groupoïde $T$-\uncertain{fidèle}), une opération de $\Sigma$ sur
$(T, \Lambda)$, et un hom.\ de groupoïdes
\[
  \Lambda \xrightarrow{\ \text{foncteur fidèle}\ } \text{groupoïde}
\]
\note{« fidèle » est souligné deux fois}
: un seul \uncertain{objet}, dont les \uncertain{autom.} sont $\pi$ ; des revêt.\
universels de $\Pi$ égaux à $\widetilde{\Pi}$ \ldots, compatible avec les
opérations de $\Sigma$. Soit $\mathcal{T} = \mathrm{Ob}\,\Lambda$ [lequel
$\mathcal{T}$ est un $\Sigma$-ens.], et pour $t, t' \in \mathcal{T}$, on a un
hom.\ \uncertain{injectif}
\[
  \underbrace{\mathrm{Isom}_{\Lambda}(t, t')}_{\pi_{t,t'}} \hookrightarrow \pi
\]
compatible avec la composition des Isom.

Ceci exprime la donnée d'un foncteur fidèle d'un groupoïde $\Lambda$ dans
$\langle \pi \rangle$ \struck{Mais on} [il faut en donner une famille de
\ill{} groupes $(\pi_t)_{t \in \mathcal{T}}$ de $\pi$ et d'ensembles \ill{}
\ill{} ces $\pi_{t'}$, $\pi_t$ \ldots] On va maintenant introduire les op.\ de

\page{24}
$\pi$ sur $\Lambda$, compatibles avec ses op.\ sur $\pi$ par automorph.\ int.
Pour ceci, on va exiger
\[
  \left\{
  \begin{array}{l}
    \forall t \in \mathcal{T} = \mathrm{Ob}\,\Lambda, \text{ la famille des }
      \pi_{t,t'} \subset \pi \ (t' \in \mathcal{T} \text{ variable}) \text{ forme une}\\
    \text{pseudo-partition de } \pi, \text{ i.e. } \forall \lambda \in \pi,\
      \exists \text{ un unique } t' \text{ tel que } \lambda \in \pi_{t,t'},\\
    \text{i.e. une unique flèche de } \Lambda \text{ au-dessus de } \lambda
      \text{ qui soit d'origine } t.
  \end{array}
  \right.
\]
\marginal{examinons \ill{} condition $\exists !$ de remontée des glissements
: origines \ill{}, \ill{} \ill{} de groupoïdes \ldots}
\note{note écrite dans la marge de gauche, à la hauteur de l'accolade}

Ce qui donne l'opération de $\pi$ sur $\mathcal{T} = \mathrm{Ob}\,\Lambda$, et
la \uncertain{vérification} de
\[
  \Lambda \simeq \langle \mathcal{T}, \pi \rangle
\]
et le foncteur $\Lambda \to \langle \pi \rangle$. \ill{} \struck{\ill{}}
\ill{} de plus que les \uncertain{stabilisateurs} $\pi_t$ sont abéliens, d'où
une famille de groupes $(T_j)_{j \in J = \mathcal{T}/\pi}$, et opération de
$\Sigma$ sur $(T, \Lambda)$, de sorte que $\Lambda \to \langle \pi \rangle$
soit compatible avec op.\ de $\Sigma$ \ldots

Ce n'est pas bien éclairant --- peut-être est-il plus agréable de

\page{25}
\note{p. 13 de l'auteur}

considérer que \struck{ensemble} $\mathcal{T}$ (discret) (qui \ill{}
l'ensemble $\pi$) \ill{} un foncteur

\begin{tikzcd}[column sep=large]
    \mathrm{Rev}(\Pi) \arrow[r, "F"] & (\mathrm{Ens}) \\
    (\pi\text{-}\mathrm{Ens}) \arrow[u, "\simeq"] \arrow[ur, bend right=20] &
  \end{tikzcd}

\note{sur la flèche verticale, un début d'étiquette biffé, puis
« $E \mapsto \widetilde{\Pi}{}^{0} \times_{\pi} E$ » ; la flèche courbe de
$(\pi\text{-}\mathrm{Ens})$ vers $(\mathrm{Ens})$ porte
« $E \mapsto \mathcal{T}^{0} \wedge_{\pi} E$ »}

[NB l'indice $0$ indique qu'on a \uncertain{remplacé} les \ill{} \ill{} \ill{}
\ill{} faire opérer $\pi$ : \ill{}) \ldots]

Foncteur \ill{} \ill{} commutant aux $\varinjlim$ quelconques. Ceci dit, on
veut que $F$ (qui opère sur $\mathrm{Rev}(\Pi)$ via ses opérations sur $\Pi$)
soit compatible avec les foncteurs $F$, i.e.\ qu'on ait des isom.
\[
  F({}^{\gamma}\Pi') \simeq F(\Pi') \qquad (\Pi' \in \ill{}\ \Pi)
\]
\note{l'indice de $\Pi' \in \ldots$ est surchargé et n'a pas été lu}
fonctoriels en $\Pi'$, et une transitivité pour $\gamma$ variable. Cela
\ill{} \ill{} \ill{}

\page{26}
que l'opération de $\pi$ sur $\mathcal{T}$ soit prolongée en une opération de
$\Sigma$ sur $\mathcal{T}$. Mais ceci ne donne pas d'interprétation
\ill{} de la $T$-\ill{}. Arrêtons \ldots

\section*{(9) Les schémas et Topos fondamentaux dans la théorie \ldots}

\note{titre de l'auteur, le 9 entouré ; la phrase s'interrompt sur des points}

\subsection*{(A) \struck{Topos \ill{}} Topos \uncertain{modulaires}}

\note{« Topos » récrit au-dessus du premier mot biffé}

\ldots\ $(\mathcal{T}_{0,\nu})_{\overline{\mathbb{Q}}}$ pour les courbes de
genre $0$ avec $\nu$ points marqués, et famille des modulaires \ill{} \ill{}
sur $\mathcal{T}_{0,\nu}$ \ldots

Le Topos $(\mathcal{T}_{0,\nu})_{\mathbb{Q}}$ pour $\nu = 0, 1, 2$
\uncertain{existe} [\uncertain{comme} topos] étendues --- car les schémas \ill{}
d'hom.\ d'une $X/S$ fibré projectif de genre $0$ et de $\nu$ sections
marquées \ill{} (disjointes). Cela tient au fait que \ill{} de lieu \ill{}
$3 - \nu > 0$ (i.e.\ $0 \leqslant \nu \leqslant 2$). \struck{\ill{}} Ce sont en
fait les topos
\[
  B_{\mathrm{PGL}(2)}, \quad B_{\mathrm{Aff}(1)}, \quad B_{\mathbb{G}_m},
\]
($B_G$ = \uncertain{classifiant} de \ill{})

\page{27}
\note{p. 14 de l'auteur}

groupes fondamentaux sont triviaux :
\[
  \pi_1(\mathcal{T}_{0,\nu}) = 0, \qquad 0 \leqslant \nu \leqslant 3 .
\]
\note{le $3$ est repassé sur un autre chiffre ; au-dessus de « triviaux »,
une addition de l'auteur, « \ill{} connue »}

Le Type $\mathcal{T}_{0,\nu}$ ($\nu \geqslant 4$) correspond au
\struck{$\to$ le topos \ill{}} schéma \struck{étale de} des syst.\ de $\nu - 3$
points distincts $(x_1, \ldots, x_{\nu-3})$ dans
\[
  \mathbb{P}^1 \smallsetminus \{0, 1, \infty\} = \mathbb{E}^1 \smallsetminus \{0, 1\} .
\]
\note{la lettre de $\mathbb{E}^1$ est un E à double trait, lu tel quel}

Son groupe fondamental est le complété profini d'un groupe de
\uncertain{tresses} $\Pi_{0,\nu}$, lequel (« groupe \ill{} ») dépend
contra-fonctoriellement de $[1, \nu]$ pour les monomorphismes. Sur
$(\mathcal{T}_{0,\nu})_{\mathbb{Q}}$ du $(\mathcal{T}_{0,\nu})_{\overline{\mathbb{Q}}}$
on fait opérer $\mathfrak{S}_\nu$, donc $\mathfrak{S}_\nu \times \Gamma$ opère
sur $\mathcal{T}_{0,\nu}$,

\note{« du $(\mathcal{T}_{0,\nu})_{\overline{\mathbb{Q}}}$ » est ajouté
au-dessus de la ligne ; le reste de la page est blanc}

\section*{Une conjecture arithmético-géométrique}

\note{titre de l'auteur, souligné, précédé d'un chiffre entouré lu 9}

\page{28}
\note{p. 15 de l'auteur}

Soient $g, \nu \geqslant 0$ deux entiers. On désigne par $M_{g,\nu}$ le topos
modulaire (sur $\mathbb{Q}$) pour les courbes (relatives) lisses propres
relativement connexes de genre $g$, munies d'un $\nu$-uple \uncertain{ordonné}
de $\nu$ sections disjointes deux à deux fibre par fibre.
\note{le $M$ est écrit en surcharge sur une première lettre, qui semble être
un $\mathcal{T}$ ; « sur $\mathbb{Q}$ » et « relatives » sont ajoutés
au-dessus de la ligne}

Sauf pour $g = 0$, $\nu = 0, 1, 2$ et $g = 1$, $\nu = 0$, ce topos est
\uncertain{associé} : \ill{} étendue algébrique, \ill{} $=$ \ill{}.

On a un hom.\ naturel \struck{extérieur}
\[
  (*) \qquad \pi_1(M_{g,\nu;\xi}) \longrightarrow
  \mathrm{Autext}_{\mathrm{lac}}^{\circ}\bigl(\pi_1(\overline{X} \smallsetminus
  \{x_1, \ldots, x_\nu\})\bigr)
\]
où $(X; x_1, \ldots, x_\nu)$ (sur $\overline{\mathbb{Q}}$) $\to$ le pt
correspondant de genre $g$ \ill{} $M_{g,\nu}(\overline{\mathbb{Q}})$ \ldots\ avec
$\nu$ pts marqués associé à $\xi$ (\uncertain{l'exposant} $\circ$ indique qu'on
se borne aux automorphismes du $\pi_1$ qui \ill{} les classes
\uncertain{d'inertie} aux $x_i$ \ldots). \uncertain{Notations} \ill{}
\[
  (**) \qquad \pi_1(\overline{M}_{g,\nu;\xi}) \longrightarrow
  \mathrm{Autext}_{\mathrm{lac}}^{\circ\circ}\bigl(\pi_1(\overline{X} \smallsetminus
  \{x_1, \ldots, x_\nu\})\bigr)
\]
où le double $\circ$ indique qu'on se borne aux \ill{} \uncertain{pour des}
\ill{} \ldots\ \uncertain{strictement} \ill{} induisant l'identité sur
$T = \mathrm{Tate}(\overline{\mathbb{Q}})$.
\[
  (***) \qquad \pi_1(M_{g,\nu;\xi}) \longrightarrow
  \mathrm{Autext}_{\mathrm{lac}}^{\circ\circ}\,
  \underbrace{\pi_1(X \smallsetminus \{x_1, \ldots, x_\nu\})}_{=\ \mathcal{T}_{g,\nu}\ (\text{Teichmüller})}
\]
(théorie top.) qui induit un

\marginal{il faut \ill{} \uncertain{plus strict}, et exiger la \ill{} aux
\ill{} --- des \ill{} \ill{} définition \ill{} $\pi_1$ ?? \ldots\ $g = 0$,
$\nu = 3$ \ldots}
\note{note écrite en diagonale dans la marge inférieure gauche ; le signe
sous l'accolade de $(***)$ est précédé d'un mot biffé illisible}

\page{29}
\uncertain{isomorphisme}, et dans ces cas on \ill{} \ill{} un \ill{} \ill{}.
On a ainsi un \ill{} :

\note{le diagramme suivant est numéroté (4) par l'auteur}

\begin{tikzcd}[column sep=small, row sep=normal, nodes={font=\scriptsize}]
    1 \arrow[r] & \pi_1(\overline{M}_{g,\nu;\xi}) \arrow[r] \arrow[d, "\wr"'] &
      \pi_1(M_{g,\nu;\xi}) \arrow[r] \arrow[d] & \Gamma \arrow[r] \arrow[drr] & 1 & \\
    1 \arrow[r] & \widehat{\mathcal{T}}_{g,\nu} \arrow[r] &
      \mathrm{Autext}_{\mathrm{lac}}^{\circ}\bigl(\pi_1(X \smallsetminus \{x_1, \ldots, x_\nu\})\bigr)
      \arrow[rrr] & & & \mathrm{Out}\ldots \arrow[d] \\
    & & & & & 1
  \end{tikzcd}

\note{la flèche verticale de gauche est marquée « $\wr$ » avec un point
d'interrogation, et porte à droite une étiquette biffée, récrite
« $\mathcal{T}_{g\nu}$ » ; le dernier terme de la seconde ligne est coupé par
le bord du feuillet et n'est lu qu'à moitié}

\marginal{sauf si $g = 0$, $\nu = 0, 1, 2$ ; $g = 1$, $\nu = 0$ ?}
\note{en haut à gauche, entouré}

\ldots\ que $\widehat{\mathcal{T}}_{g,\nu} \to
\mathrm{Autext}_{\mathrm{lac}}^{\circ}(\ill{})$ \uncertain{injective} \ill{}
dans \ldots

On \uncertain{aimerait} \ill{} que la flèche (1) (la flèche verticale centrale
de (5)) soit un \uncertain{iso}. Cela donnerait une construction de $\Gamma$
comme
\[
  \underbrace{\mathrm{Autext}_{\mathrm{lac}}^{\circ}\bigl(\pi_1(X \smallsetminus
  \{x_1, \ldots, x_\nu\})^{\wedge}\bigr)}_{\text{groupe \ill{} bien connu}}
  \big/ \widehat{\mathcal{T}}_{g,\nu}
\]

\note{« (1) » et « (5) » sont lus tels quels ; le diagramme ci-dessus est
numéroté (4)}

\emph{Ex} Pour $g = 0$, $\nu = 3$ cela donnerait l'expression déjà
envisagée de $\Gamma$ comme
\[
  \mathrm{Autext}_{\mathrm{lac}}^{\circ}\bigl(\pi_1(\mathbb{P}^1 \smallsetminus
  \{0, 1, \infty\})\bigr).
\]

\marginal{Vérifier conjecture : $\pi_1(M_{g,\nu;\xi}) \simeq$ normalisateur de
$\widehat{\mathcal{T}}_{g,\nu}$ dans
$\mathrm{Autext}_{\mathrm{lac}}^{\circ}\bigl(\pi_1(X_\xi \smallsetminus \{x_1,
\ldots\})\bigr)$ \ldots\ commutant \ill{} et structure \ill{}
\uncertain{locale} \ldots}
\marginal{p.\ ex.\ $(g, \nu)$ \ldots\ $g = 0$, $\nu = 0, 1, 2$ ; $g = 1$
\ldots\ il conviendrait d'être \ill{} l'exigence \ill{} \ldots}
\note{deux notes écrites en diagonale dans la marge de gauche ; un bloc de la
seconde, encadré, est raturé en zigzag et n'a pas été lu}

\page{30}
\note{p. 16 de l'auteur}

Si on considère le quotient $N_{g,\nu}$ de $M_{g,\nu}$ par $\mathfrak{S}_\nu$
(de sorte que $M_{g,\nu}$ est un rev.\ de groupe $\mathfrak{S}_\nu$ de
$N_{g,\nu}$) \ldots
\[
  (5) \qquad 1 \to \pi_1(\overline{M}_{g,\nu}) \longrightarrow
  \pi_1(N_{g,\nu;\xi}) \longrightarrow \mathfrak{S}_\nu \times \Gamma \to 1
\]
\[
  \pi_1(N_{g,\nu;\xi}) = \pi_1(\overline{N}_{g,\nu}, \Gamma; \xi) =
  \pi_1(M_{g,\nu}; \mathfrak{S}_\nu, \xi) = \pi_1(\overline{M}_{g,\nu},
  \mathfrak{S}_\nu \times \Gamma)
\]
\note{les trois dernières écritures sont posées sous le terme du milieu,
reliées par des signes $=$ verticaux}

structure d'extension un peu plus riche que la première ligne de (4). La
conjecture fait équivalent : la suivante :
\[
  \pi_1(N_{g,\nu;\xi}) \xrightarrow{\ \sim\ }
  \mathrm{Autext}_{\mathrm{lac}}\bigl(\pi_1(X_\xi \smallsetminus \{x_1, \ldots, x_\nu\})\bigr)
\]
on en tire
\[
  \Gamma \simeq \mathrm{Autext}_{\mathrm{lac}}\bigl(\pi_1(X_\xi \smallsetminus
  \{x_1, \ldots, x_\nu\})\bigr) \big/ \widehat{\widetilde{\mathcal{T}}}_{g,\nu}
\]
où
\[
  \widetilde{\mathcal{T}}_{g,\nu} \simeq
  \mathrm{Autext}_{\mathrm{lac}}^{+}\bigl(\pi_1(X_\xi^{\mathrm{top}} \smallsetminus
  \{x_1, \ldots, x_\nu\})\bigr).
\]
\note{l'exposant de $X_\xi$ dans la dernière formule est un petit signe
surmonté d'un trait, lu « top » sous réserve}

${}^{*}$NB Il faudrait déterminer le groupe des automorphismes
$G_{g,\nu}$ (des topos modulaires géom.\ $\overline{M}_{g,\nu}$
\struck{\ill{}}) ($G_{g,\nu} = \mathrm{Im}\,\mathfrak{S}_\nu$ ?) sauf
\ill{} $\mathfrak{S}_\nu$ \uncertain{pour} fidèlement \ill{} $g = 0$,
$\nu \leqslant 4$ \ldots)

\note{ce NB est d'une autre encre, plus claire, comme la question qui suit}

\emph{Question}
\[
  \Gamma \longrightarrow
  \mathrm{Autext}\bigl(\underbrace{\widehat{\mathcal{T}}_{g,\nu}}_{\pi_1(\overline{M}_{g,\nu})}\bigr)
  \quad
  \left\{\begin{array}{l} \text{injectif} \\ \text{bijectif} \end{array}\right\} ?
\]
(Sans doute il faut ajouter \ill{} condition \ill{}, cf.\ NB.)

\page{31}
\emph{Question} \quad $\mathrm{Autext}(\mathcal{T}_{g,\nu}) = 1$ ?

\note{la page ne porte que cette ligne}

\page{32}
\note{feuillet de listing retourné, écrit en largeur ; les deux diagrammes
sont dessinés en perspective, chaque ligne décalée vers la droite, et sont
redessinés ici sur une grille}

\begin{tikzcd}[column sep=tiny, row sep=small, nodes={font=\scriptsize}]
    1 \arrow[dr] & & 1 \arrow[dr] & & & & & & \\
    1 \arrow[r] & \pi_1(\overline{X}_\xi \smallsetminus S_\nu) \arrow[rr] \arrow[dr] & &
      \pi_1(\overline{X}_\xi \smallsetminus S_\nu) \arrow[rr] \arrow[dr] & & 1 \arrow[dr] & & & \\
    & 1 \arrow[r] & \pi_1(\overline{M}_{g,\nu+1}) \arrow[rr] \arrow[dr] & &
      \pi_1(M_{g,\nu+1}) \arrow[rr] \arrow[dr] & & \Gamma \arrow[r] \arrow[dr] & 1 & \\
    & & 1 \arrow[r] & \pi_1(\overline{M}_{g,\nu}) \arrow[rr] \arrow[dr] & &
      \pi_1(M_{g,\nu}) \arrow[rr] \arrow[dr] & & \Gamma \arrow[r] \arrow[dr] & 1 \\
    & & & & 1 & & 1 & & 1
  \end{tikzcd}

\note{le trait sur le second $X_\xi$ de la première ligne n'est pas sûr}

\begin{tikzcd}[column sep=tiny, row sep=small, nodes={font=\scriptsize}]
    1 \arrow[dr] & & 1 \arrow[dr] & & & & & & \\
    1 \arrow[r] & \widehat{\pi}_{g,\nu} \arrow[rr, "\sim"] \arrow[dr] & &
      \widehat{\pi}_{g,\nu} \arrow[rr] \arrow[dr] & & 1 \arrow[dr] & & & \\
    & 1 \arrow[r] & \widehat{\mathcal{T}}_{g,\nu+1} \arrow[rr] \arrow[dr] & &
      \mathrm{Aut}_{\mathrm{lac}}^{\circ}(\widehat{\pi}_{g,\nu})^{!} \arrow[rr] \arrow[dr] & &
      \Gamma \arrow[r] \arrow[dr, no head, "=" description] & 1 & \Sigma_{\nu+1} \\
    & & 1 \arrow[r] & \widehat{\mathcal{T}}_{g,\nu} \arrow[rr] \arrow[dr] & &
      \mathrm{Autext}_{\mathrm{lac}}^{\circ}(\widehat{\pi}_{g,\nu})^{!} \arrow[rr] \arrow[dr] & &
      \Gamma \arrow[r] & 1 \quad \Sigma_\nu \\
    & & & & 1 & & 1 & &
  \end{tikzcd}

\note{sous le premier $\widehat{\pi}_{g,\nu}$, une écriture biffée
(« $\pi_1(\overline{X}) \ldots$ » ?) ; sous
$\mathrm{Autext}_{\mathrm{lac}}^{\circ}(\widehat{\pi}_{g,\nu})^{!} \to \Gamma$,
un mot souligné d'une accolade, \uncertain{tresses} ; la liaison entre les
deux $\Gamma$ est un double trait. Les étiquettes $\Sigma_{\nu+1}$ et
$\Sigma_\nu$ sont écrites loin à droite des deux lignes}

\[
  M'_{g,\nu}(k) \ni \xi, \qquad \mathcal{T}'_{g,\nu} \ni G_\xi
\]
\note{en bas à gauche du feuillet}

\[
  \widehat{\widetilde{\Sigma}}_{g,\nu} = \overline{\Sigma}_{g,\nu}
  \qquad\qquad \Sigma_{g,\nu}
\]
\note{en bas, au milieu : le premier terme est précédé et surmonté de signes raturés, illisibles, et
le $=$ est vertical, $\overline{\Sigma}_{g,\nu}$ étant écrit sous
$\widehat{\widetilde{\Sigma}}_{g,\nu}$}

\end{document}
